\section{The Virtue of Religion and the Needs of the Worshiper}

If worship belongs to justice without equivalence, what kind of virtue is religion itself, what acts belong properly to it, and why does an infinitely self-sufficient God require outward bodily reverence from creatures of dust and spirit? St.\ Thomas Aquinas answers these questions across five tightly linked questions of the \emph{Secunda Secundae} (\emph{ST} II-II, qq.\ 81--85), and his argument converges strikingly with the early Church's own defense of Christian worship in St.\ Irenaeus of Lyons and St.\ Augustine.

\subsection{Chief of the Moral Virtues, Commanded by Charity}

Aquinas opens his treatment of religion by tracing the three classical etymologies inherited from Cicero and St.\ Augustine: \emph{relegere}, to read over and ponder diligently what belongs to divine worship; \emph{religere}, Augustine's wordplay in \emph{City of God} X.3, to choose God again after losing him through negligence; and \emph{religare}, Augustine's word in \emph{De vera religione}, to bind ourselves to the one Almighty God.\endnote{Aquinas, \emph{ST} II-II, q.\ 81, a.\ 1, corpus; cf.\ Augustine, \emph{De civitate Dei} X.3 (CSEL 40/1: \emph{Hunc eligentes uel potius religentes (amiseramus enim neglegentes)}).} All three converge on one relation: God is the unfailing principle to whom we are bound, the last end toward whom choice must turn, and the good we recover through faith when sin has severed us.

Because God possesses a singular, infinite excellence as first principle of the creation and government of all things, the honor due to him is singular as well, known in Greek as \emph{latria} (the supreme service due to God alone). Religion is therefore one special virtue, distinct both from the virtues by which we honor parents or rulers and from general uprightness (\emph{ST} II-II, q.\ 81, aa.\ 3--4).

Yet Aquinas draws a decisive distinction that keeps worship from swallowing theology or displacing grace: \emph{religion is not a theological virtue}.\endnote{Aquinas, \emph{ST} II-II, q.\ 81, aa.\ 5--6.} A theological virtue---faith, hope, or charity---has God himself as its direct, immediate object: by faith we believe God himself, by hope we lean directly on divine help for beatitude, and by charity we love God in spiritual friendship. Religion, by contrast, has as its immediate matter and object the \emph{acts of worship} we offer---devotion, prayer, adoration, sacrifice, vows---while God is its \emph{end}, to whom those acts are directed. Religion is therefore a moral virtue, not a theological one.

At the same time, religion stands at the summit of all moral virtues (\emph{ST} II-II, q.\ 81, a.\ 6). Whatever is ordered to an end draws its goodness from proximity to that end; because religion's acts are directed immediately to the honor of God, it surpasses fortitude, temperance, and ordinary social justice among the moral virtues. Yet because the virtue concerned with the end moves and commands the virtues concerned with the means, the theological virtues command religion from above: ``the theological virtues, faith, hope and charity have an act in reference to God as their proper object: wherefore, by their command, they cause the act of religion, which performs certain deeds directed to God'' (\emph{ST} II-II, q.\ 81, a.\ 5, ad 1).

From the opposite direction, religion itself exercises a commanding office over the rest of human life. Aquinas resolves the famous verse of St.\ James---``Religion clean and undefiled before God and the Father is this: to visit the fatherless and widows in their tribulation and to keep one's self unspotted from this world'' (Jas 1:27)---by distinguishing \emph{elicited} acts from \emph{commanded} acts.\endnote{Aquinas, \emph{ST} II-II, q.\ 81, a.\ 1, ad 1; cf.\ II-II, q.\ 81, a.\ 8 on religion and sanctity (\emph{sanctitas}).} Religion's immediate, elicited acts direct man to God alone: devotion, prayer, adoration, sacrifice. Its commanded acts belong immediately to other virtues---mercy visits the widow and orphan, temperance keeps one unspotted from the world---yet religion directs them to the honor of God. Viewed specifically in its worshipful acts, this virtue is called religion; viewed as referring all virtuous deeds to God and keeping the mind pure and steadfast in him, it is called \emph{sanctity}. Worship at the altar and uprightness in the world are neither identical nor separable: they share one root, ordered from charity down through religion into every virtue.

\subsection{Interior and Exterior Acts: Devotion, Prayer, Adoration}

Within religion's own elicited acts, Aquinas establishes a clear hierarchy: interior acts hold the first place essentially, while exterior bodily acts are secondary and ordered to the interior (\emph{ST} II-II, q.\ 81, a.\ 7; q.\ 84, a.\ 2). That ordering never demotes the body to irrelevance; rather, it explains why embodied worship is indispensable for a creature who is neither an angel nor a disembodied mind.

The primary interior act of religion is \emph{devotion}.\endnote{Aquinas, \emph{ST} II-II, q.\ 82, aa.\ 1--3.} Contrary to modern usage, which often equates devotion with tender emotion, Aquinas defines it from its Latin root \emph{devovere} (to vow or dedicate) as a firm act of the will: ``the will to give oneself readily to things concerning the service of God'' (\emph{ST} II-II, q.\ 82, a.\ 1). Charity causes devotion, because love makes a friend eager to serve; yet devotion in turn nourishes charity: ``Even so all friendship is safeguarded and increased by the practice and consideration of friendly deeds'' (\emph{ST} II-II, q.\ 82, a.\ 2, ad 2). And how is devotion kindled within us? While God is its chief and extrinsic cause, its intrinsic cause on our side is meditation on two paired truths: God's goodness and loving-kindness, which awakens love, and our own frailty and need, which shuts out presumption (\emph{ST} II-II, q.\ 82, a.\ 3). Because the human mind is weak and learns through sensible things, Aquinas notes that the chief guiding hand to devotion is meditation on the sacred humanity of Christ, echoing the Christmas Preface: ``that through knowing God visibly, we may be caught up to the love of things invisible'' (\emph{ST} II-II, q.\ 82, a.\ 3, ad 2).

Next to devotion stands \emph{prayer}.\endnote{Aquinas, \emph{ST} II-II, q.\ 83, aa.\ 1--3, 12.} Drawing on Cassiodorus and St.\ John of Damascus, Aquinas calls prayer \emph{oris ratio} (spoken reason) and the raising of the mind to God to ask becoming things. Crucially, prayer neither informs God of needs he does not know nor alters the eternal decree of providence. Divine providence disposes not only which effects will come about, but also through which secondary causes they will come about; among those causes are human prayers. We pray, Aquinas writes, ``not that we may change the Divine disposition, but that we may impetrate that which God has disposed to be fulfilled by our prayers, in other words `that by asking, men may deserve to receive what Almighty God from eternity has disposed to give,' as Gregory says'' (\emph{ST} II-II, q.\ 83, a.\ 2). By asking, the creature confesses openly that God is the author of every good and learns confidence in turning to him.

Moreover, prayer naturally passes into human voice (\emph{ST} II-II, q.\ 83, a.\ 12). Common prayer offered by the Church's ministers representing the whole body must be vocal so that the assembled faithful know and join what is offered. Even private prayer rightly employs words for three reasons: to arouse interior devotion through sensible signs, to pay our debt with the body as well as the mind so that we serve God with all we have received from him, and out of the natural overflow of a glad heart onto the tongue.

The same anthropology governs \emph{adoration}.\endnote{Aquinas, \emph{ST} II-II, q.\ 84, aa.\ 1--2.} Because human nature is composite---intellectual and sensible at once---we owe God both spiritual adoration in the mind and bodily adoration in posture and gesture. Genuflecting signifies our weakness before God; prostrating ourselves confesses that of ourselves we are nothing (\emph{ST} II-II, q.\ 84, a.\ 2, ad 2). Far from violating Christ's word to the Samaritan woman that true adorers adore the Father ``in spirit and in truth'' (John 4:23--24), bodily adoration is itself done in spirit whenever it springs from and feeds interior reverence (\emph{ST} II-II, q.\ 84, a.\ 2, ad 1). At the same time, Aquinas keeps \emph{latria} (adoration of supreme excellence, reserved to God alone) strictly distinct from \emph{dulia} (reverence paid to saints or honored creatures), noting with Augustine that while many signs of honor are shared analogously with men, sacrifice belongs to God alone (\emph{ST} II-II, q.\ 84, a.\ 1, ad 1).

\subsection{The Anti-Magic Principle: God Needs Nothing}

Behind Aquinas's entire account of interior and exterior worship stands a principle that cleanly separates Catholic worship from pagan magic, transactional appeasement, and spiritual flattery. Magic treats ritual as a lever to bend or feed a needy deity. Catholic theology begins from the opposite premise: God lacks nothing, gains nothing from our praise, and suffers no loss from our refusal. Aquinas states the rule plainly:

\begin{quote}
We pay God honor and reverence, not for His sake (because He is of Himself full of glory to which no creature can add anything), but for our own sake, because by the very fact that we revere and honor God, our mind is subjected to Him; wherein its perfection consists, since a thing is perfected by being subjected to its superior, for instance the body is perfected by being quickened by the soul, and the air by being enlightened by the sun.\endnote{Aquinas, \emph{ST} II-II, q.\ 81, a.\ 7, corpus (English Dominican translation, Project Gutenberg eBook 18755).}
\end{quote}

The Psalter had already mocked every transactional fantasy of divine hunger: ``If I should be hungry, I would not tell thee: for the world is mine, and the fulness thereof. Shall I eat the flesh of bullocks? or shall I drink the blood of goats? Offer to God the sacrifice of praise: and pay thy vows to the most High'' (Ps 49:12--14; Heb.\ 50:12--14). External offerings are never presented as though God stood in need of them, but as signs of the inward submission by which the human mind is ordered rightly to its Creator (\emph{ST} II-II, q.\ 81, a.\ 7, ad 2).

More than a millennium before Aquinas, St.\ Irenaeus of Lyons built Book IV, chapters 14 and 17--18 of his \emph{Adversus Haereses} (\emph{Against Heresies}, c.\ 180) around this very truth, using an image that arrests the mind immediately:

\begin{quote}
Nor did He stand in need of our service when He ordered us to follow Him; but He thus bestowed salvation upon ourselves. For to follow the Saviour is to be a partaker of salvation, and to follow light is to receive light. But those who are in light do not themselves illumine the light, but are illumined and revealed by it: they do certainly contribute nothing to it, but, receiving the benefit, they are illumined by the light. Thus, also, service [rendered] to God does indeed profit God nothing, nor has God need of human obedience; but He grants to those who follow and serve Him life and incorruption and eternal glory\dots\ For, as much as God is in want of nothing, so much does man stand in need of fellowship with God. For this is the glory of man, to continue and remain permanently in God's service.\endnote{St.\ Irenaeus of Lyons, \emph{Adversus Haereses} IV.14.1 (Roberts--Rambaut translation, \emph{Ante-Nicene Fathers}, vol.\ 1).}
\end{quote}

Irenaeus applies the same logic to the whole Mosaic pedagogy of tabernacle, temple, and Levitical sacrifices (\emph{Adversus Haereses} IV.14.3; IV.17.1). God had in himself ``every perfume of sweet-smelling savours, even before Moses existed.'' He enjoined outward observances not because he hungered for blood or incense, nor in petty vanity, but out of compassion for a people prone to idols, teaching them through earthly types to fear God, to persevere in his service, and to draw near to him by obedience, contrition, and righteousness.

St.\ Augustine compresses the same principle into two sentences in \emph{City of God} X.5:

\begin{quote}
We must believe, then, that God has no need, not only of cattle, or any other earthly and material thing, but even of man's righteousness, and that whatever right worship is paid to God profits not Him, but man. For no man would say he did a benefit to a fountain by drinking, or to the light by seeing.\endnote{St.\ Augustine, \emph{De civitate Dei} X.5 (Dods 1871; CSEL 40/1: \emph{Neque enim fonti se quisquam dixerit consuluisse, si biberit; aut luci, si uiderit}).}
\end{quote}

Refusing to drink leaves the traveler dead of thirst, and shutting one's eyes leaves him in the dark. Worship is owed to God in truth because he is the Source; it benefits not God, who is complete in eternal beatitude, but the creature whose life and splendor consist in being turned toward the light.
